
That downstream performance degrades during spectral expansion confirms these are collapse indicators, not quality measures; SLV and perplexity are complementary.
Our analysis reveals a reproducible three-phase collapse trajectory, identifies SLV as a cheap spectral indicator, and uncovers a prevention--detectability tradeoff: real-data mixing reverses the spectral trajectory and the controller signal vanishes at $r \geq 0.3$, so rate-of-change-based reactive controllers cannot function where prevention succeeds.
The trajectory map localizes $r_{\text{crit}} \in (0.25, 0.3]$; setting $r > r_{\text{crit}}$ prevents collapse without per-generation monitoring.
